% Responsable: P3
\section{Léxico, concordancia y visualización}
\label{sec:lexer_concordancia_arboles}

\subsection{El lexer}

El lexer (\texttt{spanish/lexer.py}) es una adaptación del de la Fase~1. Recorre la oración
carácter por carácter, agrupa las letras en palabras y, como antes, reporta los errores y sigue
en lugar de detenerse. Los cambios respecto a la Fase~1 son:
\begin{itemize}
  \item \textbf{Alfabeto.} Las 26 letras latinas más \emph{á, é, í, ó, ú, ü, ñ} y sus mayúsculas.
    No se usa \texttt{str.isalpha()}, que también aceptaría letras de otros alfabetos
    (cirílicas, griegas o \emph{ç}); esas se reportan como caracteres inválidos.
  \item \textbf{Normalización.} Algunos sistemas guardan \emph{é} como \emph{e} más un acento
    combinado (forma NFD). El texto se normaliza a NFC para que cada letra acentuada sea un solo
    carácter y coincida con el léxico.
  \item \textbf{Puntuación.} Se ignora un único signo final (\texttt{.}, \texttt{!} o
    \texttt{?}); cualquier otro signo es un error léxico.
  \item \textbf{La categoría viene del léxico, no de la forma.} En la Fase~1 una expresión regular
    decidía si un token era un número o un identificador. En español la misma forma puede ser
    sustantivo, verbo o adjetivo, así que cada palabra recibe todas sus lecturas del léxico y el
    parser elige (sección~\ref{sec:parser}).
  \item \textbf{Palabras desconocidas.} Una palabra que no está en el léxico no es un error
    léxico: recibe la categoría \texttt{X} (fuera de la gramática) y el parser rechaza la oración
    con error sintáctico. Así la salida conserva la etiqueta de cada palabra.
\end{itemize}

\subsection{El léxico}

El léxico (\texttt{data/lexicon.json}) tiene 221 palabras. Cada una lleva su categoría y sus
rasgos:
\begin{lstlisting}
"niña":  {"cat": "Noun", "gen": "f", "num": "sg"},
"mató":  {"cat": "Verb", "lema": "matar", "tiempo": "pret", "pers": 3, "num": "sg"},
"bajo":  [{"cat": "Adj", "gen": "m", "num": "sg"}, {"cat": "Prep"}],
\end{lstlisting}
Una palabra ambigua lleva una lista de entradas, como \emph{bajo}.

\begin{table}[h]
  \centering
  \small
  \begin{tabular}{lrl}
    \toprule
    Categoría & Lecturas & Rasgos \\
    \midrule
    \textit{Prn} (pronombre) & 8 & género, número, persona \\
    \textit{Det} (determinante) & 8 & género, número \\
    \textit{Pos} (posesivo) & 6 & número \\
    \textit{Noun} (sustantivo) & 52 & género, número \\
    \textit{Adj} (adjetivo) & 58 & género (salvo invariables), número \\
    \textit{Intens} (intensificador) & 4 & --- \\
    \textit{Adv} (adverbio) & 9 & --- \\
    \textit{Prep} (preposición) & 11 & --- \\
    \textit{Verb} (verbo) & 70 & tiempo, persona, número \\
    \bottomrule
  \end{tabular}
  \caption{Contenido del léxico. Los verbos son 7 (\emph{buscar, comer, comprar, dormir, leer,
    matar, ver}) en presente y pretérito, en cinco personas (todas menos la segunda del plural,
    porque el léxico no incluye \emph{vosotros}). Hay más lecturas que palabras porque cinco
    palabras son ambiguas.}
  \label{tab:lexico}
\end{table}

El vocabulario se eligió para cubrir el conjunto de prueba: sustantivos de los dos géneros en
singular y plural, adjetivos variables (\emph{rojo, roja, rojos, rojas}) e invariables
(\emph{grande, gris}), los tres tipos de sujeto y las tres personas gramaticales. Los dos tiempos verbales se
resuelven en el léxico con el rasgo \texttt{tiempo}, sin reglas nuevas en la gramática.

\subsection{Verificación de la concordancia}

El analizador de concordancia (\texttt{spanish/agreement.py}) es la versión de la Fase~2 del
analizador semántico. Recorre el árbol que construyó el parser y aplica dos reglas:
\begin{enumerate}
  \item \textbf{Dentro de cada sintagma nominal}, el determinante o posesivo y el adjetivo deben
    coincidir en género y número con el sustantivo. Se revisan todos los sintagmas nominales del
    árbol: el sujeto, el objeto y los que están dentro de frases preposicionales.
  \item \textbf{Sujeto-verbo}: el núcleo del sujeto y el verbo deben coincidir en número y
    persona. Los sustantivos son de tercera persona; los pronombres llevan la suya.
\end{enumerate}
Un rasgo ausente concuerda con cualquier valor. Así se tratan los adjetivos invariables
(\emph{grande} no tiene género) y los posesivos, que solo marcan número (\emph{su}, \emph{sus}).
Como la verificación se hace sobre el árbol y no sobre palabras vecinas, funciona a distancia:
en \emph{El perro de los niños come} el verbo concuerda con \emph{perro}, el núcleo del sujeto,
y no con \emph{niños}, la palabra más cercana. Por eso se acepta, mientras que
\emph{*El perro de los niños comen} se rechaza.

Cada error se reporta en la palabra que rompe la concordancia y dice qué rasgo falla, por
ejemplo: \emph{el adjetivo `grises' no concuerda en número con `elefante'}.

\subsection{Visualización}

\texttt{spanish/trees.py} dibuja los árboles de dos formas: en la terminal, con
\texttt{nltk} (opción \texttt{--arbol} del pipeline), y como imagen PNG, SVG o PDF con Graphviz
(opción \texttt{--png}). En la imagen cada categoría tiene su color y las palabras que rompen la
concordancia aparecen en rojo, como en la figura~\ref{fig:arbol_concordancia}.

\begin{figure}[h]
  \centering
  \begin{forest} sintaxis
    [S [NP [NPbase [Det [\textcolor{red}{la}]]
                   [Nom [Noun [elefante]] [AdjP [Adj [\textcolor{red}{grises}]]]]]]
       [VP [Verb [duerme]]]]
  \end{forest}
  \caption{\emph{*La elefante grises duerme}: la gramática la deriva y el analizador de
    concordancia marca las dos palabras que no concuerdan con \emph{elefante}.}
  \label{fig:arbol_concordancia}
\end{figure}
